\documentclass[12pt,a4paper]{exam}
\usepackage{geometry}
\usepackage{fancyvrb}
\usepackage{pdfsync}

\geometry{
  %includeheadfoot,
  margin=2.54cm
}

\title{Exercises --- Scala (Week Three)}
\author{Introduction to the language}
\date{Spring term 2017}

\begin{document}

\maketitle

\noindent 
You need to have installed the Scala distribution before commencing these exercises.
Don't forget the documentation for the distribution as that will come in very useful.

This exercise sheet provides additional examination of the basic 
imperative and object-oriented capabilities of the language.

\begin{questions}
\question
\begin{parts}
\part
Verify your Java version by typing \verb!java –version! in a shell (command window). 
The version must be 1.6 or greater for Scala to work.
\part
Verify your Scala version by typing \verb!scala! in a shell (this starts the REPL). 
The version must be 2.10 or greater.
\part
Quit the REPL by typing \verb!:quit!.
\end{parts}

\begin{solution}
	\ldots	
\end{solution}

\question

\begin{parts}
\part Create a \emph{value} identifier and store (and print) the value \verb!17!.
\begin{solution}
\begin{Verbatim}
val v1 = 17
println(v1)
\end{Verbatim}	
\end{solution}

\part Using the value you just stored (\verb!17!), try to change it to \verb!20!. 
What happened?
\begin{solution}
\begin{Verbatim}
/*{oldDescription}
Using the value you just stored (17), try to change it to 20. 
What happened?
{oldDescription}*/

// this won't work
// v2 = 20
println("You can't change a value, once set")
\end{Verbatim}	
\end{solution}

\part Store (and print) the value \verb!ABC1234!.
\begin{solution}
\begin{Verbatim}
val s = "ABC1234"
println(s)	
\end{Verbatim}	
\end{solution}

\part Using the value you just stored (\verb!ABC1234!), try to change it to
\verb!DEF1234!. What happened?
\begin{solution}
\begin{Verbatim}
// this won't work
// s = "DEF1234"
println("You can't change a value, once set")
\end{Verbatim}	
\end{solution}

\part Store the value \verb!15.56!. Print it.
\begin{solution}
\begin{Verbatim}
val d1 = 15.56D
println(d1)	
\end{Verbatim}	
\end{solution}

\end{parts}

\question

\begin{parts}
\part
Write an expression that evaluates to \verb!true! if the sky is \verb!“sunny”! and
the temperature is more than 80 degrees.
\begin{solution}
\begin{Verbatim}
val sky = "sunny"
var temperature = 87

val sunnyAndWarm = sky == "sunny" && temperature > 80
println(sunnyAndWarm)	
\end{Verbatim}	
\end{solution}

\part 
Write an expression that evaluates to \verb!true! if the sky is either 
\verb!“sunny”! or \verb!“partly cloudy”! and the temperature is more than 80 degrees.
\begin{solution}
\begin{Verbatim}
val sky = "sunny"
var temperature = 87
val somewhatSunnyAndWarm =
  (sky == "sunny" ||
  sky == "party cloudy") &&
  temperature > 80

println(somewhatSunnyAndWarm)	
\end{Verbatim}	
\end{solution}

\part 
Write an expression that evaluates to \verb!true! if the sky is either 
\verb!“sunny”! or \verb!“partly cloudy”! and the temperature is either more than 80 degrees 
or less than 20 degrees.
\begin{solution}
\begin{Verbatim}
val sky = "sunny"
val temperature = 10

val sunnyAndExtreme =
  (sky == "sunny" ||
  sky == "party cloudy") &&
  (temperature > 80 ||
  temperature < 20)

println(sunnyAndExtreme)	
\end{Verbatim}	
\end{solution}

\part Convert Fahrenheit to Celsius. \\
Hint: first subtract 32, then multiply by 5/9. If you get 0, check to make sure you didn’t do integer maths.
\begin{solution}
\begin{Verbatim}
val faren = 80.0
val celsius = (faren-32.0) * (5.0/9.0)
println(celsius)	
\end{Verbatim}	
\end{solution}

\part Convert Celsius to Fahrenheit. \\
Hint: first multiply by 9/5, then add 32. Use this to check your solution for the previous exercise.
\begin{solution}
\begin{Verbatim}
val c = 26.67
val f = c*(9.0/5) + 32
println(f)	
\end{Verbatim}	
\end{solution}

\end{parts}

\question

\begin{parts}
\part
Create a method \verb!getSquare! that takes an \verb!Int! argument and 
returns its square. Print your answer. Test using the following code.
\begin{Verbatim}
val a = getSquare(3)
assert(/* fill this in */)

val b = getSquare(6) 
assert(/* fill this in */) 

val c = getSquare(5) 
assert(/* fill this in */)
\end{Verbatim}
\begin{solution}
\begin{Verbatim}
def getSquare(num: Int):Int = {
  num*num
}

val a = getSquare(3)
println(a)
assert(a == 9, "Expected 9, Got " + a)
val b = getSquare(6)
println(b)
assert(b == 36, "Expected 36, Got " + b)
val c = getSquare(5)
println(c)
assert(c == 25, "Expected 25, Got " + c)	
\end{Verbatim}	
\end{solution}

\part
Create a method \texttt{isArg1GreaterThanArg2} that takes two \texttt{Double} arguments. 
Return \texttt{true} if the first argument is greater than the second. 
Return \texttt{false} otherwise. 
Print your answer. Satisfy the following:
\begin{Verbatim}
val t1 = isArg1GreaterThanArg2(4.1, 4.12) 
assert(/* fill this in */)

val t2 = isArg1GreaterThanArg2(2.1, 1.2) 
assert(/* fill this in */)
\end{Verbatim}
\begin{solution}
\begin{Verbatim}
def isArg1GreaterThanArg2(num1:Double,
	num2:Double) = {
  num1 > num2
}

val t1 =
  isArg1GreaterThanArg2(4.1, 4.12)
println(t1)
assert(false == t1,
	"Expected false, Got " + t1)
val t2 = isArg1GreaterThanArg2(2.1, 1.2)
println(t2)
assert(true == t2,
	"Expected true, Got " + t2)	
\end{Verbatim}	
\end{solution}

\part
Create a method \texttt{manyTimesString} that takes a \texttt{String} and an \texttt{Int} 
as arguments and returns the \texttt{String} duplicated that many times. 
Print your answer. Satisfy the following:
\begin{Verbatim}
val m1 = manyTimesString("abc", 3)
assert("abcabcabc" == m1, "Your message here")

val m2 = manyTimesString("123", 2) 
assert("123123" == m2, "Your message here")
\end{Verbatim}
\begin{solution}
\begin{Verbatim}
def manyTimesString(str: String, num: Int) = str*num

val m1 = manyTimesString("abc", 3)
println(m1)
assert("abcabcabc" == m1, "Expected abcabcabc, Got " + m1)
val m2 = manyTimesString("123", 2)
println(m2)
assert("123123" == m2, "Expected 123123, Got " + m2)	
\end{Verbatim}	
\end{solution}

\end{parts}

\question

\begin{parts}
\part
Create a \texttt{Range} object and print the \texttt{step} value. 
Create a second \texttt{Range} object with a \texttt{step} value of \verb!2! 
and then print the \texttt{step} value. 
What's different?
\begin{solution}
\begin{Verbatim}
val r1 = Range(0, 10)
println(r1.step)

val r2 = Range(0, 10, 2)
println(r2.step)

// BONUS
val r3 = r1.by(4)
println(r3.step)	
\end{Verbatim}	
\end{solution}

\part
Create a \texttt{String} object \verb!s1! (as a \verb!var!) initialised to \verb!"Sally"!. 
Create a second \verb!String! object \verb!s2! (as a \verb!var!) initialised to \verb!"Sally"!. 
Use \verb!s1.equals(s2)! to determine if the two \verb!String!s are equivalent. 
If they are, print \verb!“s1 and s2 are equal”!, otherwise print \verb!“s1 and s2 are not equal”!.
\begin{solution}
\begin{Verbatim}
var s1 = "Sally"
var s2 = "Sally"
if (s1.equals(s2)) {
  println("s1 and s2 are equal")
}
else {
  println("s1 and s2 are not equal")
}	
\end{Verbatim}	
\end{solution}
\end{parts}

\question
\begin{parts}
\part
Create classes for \texttt{Hippo}, \texttt{Lion}, \texttt{Tiger}, \texttt{Monkey}, and \texttt{Giraffe}, 
then create an instance of each one of those classes. Display the objects. 
Do you see five different ugly-looking (but unique) strings? Count and inspect them.
\begin{solution}
\begin{Verbatim}
class Hippo
class Lion
class Tiger
class Monkey
class Giraffe

val hippo = new Hippo
val lion = new Lion
val tiger = new Tiger
val monkey = new Monkey
val giraffe = new Giraffe

println(hippo)
println(giraffe)
println(tiger)
println(monkey)
println(giraffe)	
\end{Verbatim}	
\end{solution}

\part
Create a second instance of \texttt{Lion} and two more \texttt{Giraffe}s. 
Print those objects. How do they differ from the original objects that you created?
\begin{solution}
\begin{Verbatim}
class Lion
class Giraffe

val lion2 = new Lion
val giraffe2 = new Giraffe
val giraffe3 = new Giraffe

println(lion2)
println(giraffe2)
println(giraffe3)	
\end{Verbatim}	
\end{solution}

\end{parts}

\question

\begin{parts}
\part
Create a \texttt{Sailboat} class with methods to raise and lower the sails, 
printing \texttt{“Sails raised”}, and \texttt{“Sails lowered”}, respectively. 

Create a \texttt{Motorboat} class with methods to start and stop the motor, 
returning \texttt{“Motor on”}, and \texttt{“Motor off”}, respectively. 
Create an object (instance) of the \texttt{Sailboat} class. 
Use \texttt{assert} for verification:
\begin{Verbatim}
val sailboat = new Sailboat
val r1 = sailboat.raise()
assert(r1 == "Sails raised", "Expected Sails raised, Got " + r1)

val r2 = sailboat.lower()
assert(r2 == "Sails lowered", "Expected Sails lowered, Got " + r2) 

val motorboat = new Motorboat
val s1 = motorboat.on()
assert(s1 == "Motor on", "Expected Motor on, Got " + s1)

val s2 = motorboat.off()
assert(s2 == "Motor off", "Expected Motor off, Got " + s2)
\end{Verbatim}

\begin{solution}
\begin{Verbatim}
class Sailboat {
	def raise():String = { "Sails raised" }
	def lower():String = { "Sails lowered" }
}

class Motorboat {
	def start():String = { "Motor on" }
	def stop():String = { "Motor off" }
}

val sailboat = new Sailboat
val r1 = sailboat.raise()
assert(r1 == "Sails raised",
  "Expected Sails raised, Got " +
  r1)
println(r1)

val r2 = sailboat.lower()
assert(r2 == "Sails lowered",
  "Expected Sails lowered, Got " +
  r2)
println(r2)

val motorboat = new Motorboat
val s1 = motorboat.start()
assert(s1 == "Motor on",
  "Expected Motor on, Got " +
   s1)
val s2 = motorboat.stop()
assert(s2 == "Motor off",
  "Expected Motor off, Got " +
  s2)
println(s2)	
\end{Verbatim}	
\end{solution}

\part
Create a new class \texttt{Flare}. 
Define a \texttt{light} method in the \texttt{Flare} class. 
Satisfy the following:
\begin{Verbatim}
val flare = new Flare
val f1 = flare.light
assert(f1 == "Flare used!", "Expected Flare used!, Got " + f1)
\end{Verbatim}

\begin{solution}
\begin{Verbatim}
class Flare {
	def light(): String = { "Flare used!"}
}

val flare = new Flare
val f1 = flare.light
assert(f1 == "Flare used!",
  "Expected Flare used!, Got " +
  f1)
println(f1)	
\end{Verbatim}	
\end{solution}

\part
In each of the \texttt{Sailboat} and \texttt{Motorboat} classes, 
add a method \texttt{signal} that creates a \texttt{Flare} object and calls the \texttt{light} 
method on the \texttt{Flare}. Satisfy the following:
\begin{Verbatim}
val sailboat2 = new Sailboat2
val signal = sailboat2.signal() 
assert(signal == "Flare used!", "Expected Flare used! Got " + signal) 

val motorboat2 = new Motorboat2
val flare2 = motorboat2.signal() 
assert(flare2 == "Flare used!", "Expected Flare used!, Got " + flare2)
\end{Verbatim}

\begin{solution}
\begin{Verbatim}
class Flare {
	def light(): String = { "Flare used!"}
}

class Sailboat2 {
	def raise():String = { "Sails raised" }
	def lower():String = { "Sails lowered" }
	def signal(): String = new Flare().light()
}

class Motorboat2 {
	def start():String = { "Motor on" }
	def stop():String = { "Motor off" }
	def signal(): String = new Flare().light()
}

val sailboat2 = new Sailboat2
val signal = sailboat2.signal()
assert(signal == "Flare used!",
  "Expected Flare used! Got " +
  signal)
println(signal)

val motorboat2 = new Motorboat2
val flare2 = motorboat2.signal()
assert(flare2 == "Flare used!",
  "Expected Flare used!, Got " +
  flare2)
println(flare2)	
\end{Verbatim}	
\end{solution}

\end{parts}

\question

Given the following code:
\begin{Verbatim}
class Cup {
var percentFull = 0
val max = 100
def add(increase:Int):Int = {
    percentFull += increase
    if(percentFull > max) {
      percentFull = max
    }
    percentFull // Return this value
  }
}
\end{Verbatim}

\begin{parts}
\part
What happens in \texttt{Cup}’s \texttt{add} method if \texttt{increase} is a 
negative value? Is any additional code necessary to satisfy the following tests:
\begin{Verbatim}
val cup = new Cup
cup.add(45) is 45
cup.add(-15) is 30
cup.add(-50) is -20
\end{Verbatim}
Remember to include the \texttt{AtomicTest} class which you will find it under the 
\verb!atomicscala! folder in the repo. You will need to import the methods 
from the class using the following statement:
\begin{Verbatim}
import com.atomicscala.AtomicTest._
\end{Verbatim}
\begin{solution}
\begin{Verbatim}
class Cup {
  var percentFull = 0
  val max = 100
  def add(increase:Int):Int = {
    percentFull += increase
    if(percentFull > max) percentFull = max
    percentFull // Return this value
  }
}

val cup = new Cup
cup.add(45) is 45
cup.add(-15) is 30
cup.add(-50) is -20
\end{Verbatim}	
\end{solution}

\part
Add code to handle negative values to ensure that the total never goes below 0. 
Satisfy the following tests:
\begin{Verbatim}
val cup = new Cup
cup.add(45) is 45
cup.add(-55) is 0
cup.add(10) is 10
cup.add(-9) is 1
cup.add(-2) is 0
\end{Verbatim}

\begin{solution}
\begin{Verbatim}
class Cup {
  var percentFull = 0
  val max = 100
  def add(increase:Int):Int = {
    percentFull += increase
    if (percentFull > max) percentFull = max
    if (percentFull < 0) percentFull = 0
    percentFull // Return this value
  }
}
val cup = new Cup
cup.add(45) is 45
cup.add(-55) is 0
cup.add(10) is 10
cup.add(-9) is 1
cup.add(-2) is 0	
\end{Verbatim}	
\end{solution}

\part
Can you set \texttt{percentFull} from outside the class? 
Try it, like this: 
\begin{Verbatim}
cup.percentFull = 56
cup.percentFull is 56
\end{Verbatim}

\part
Write methods that allow you to both set and get the value of \texttt{percentFull}. 
Satisfy the following:
\begin{Verbatim}
val cup = new Cup
cup.set(56)
cup.get() is 56
\end{Verbatim}

\begin{solution}
\begin{Verbatim}
class Cup4 {
  var percentFull = 0
  val max = 100
  def add(increase:Int):Int = {
    percentFull += increase
    if(percentFull > max) percentFull = max
      percentFull // Return this value
  }
  def set(x: Int) { percentFull = x }
  def get(): Int = percentFull
}

val cup4 = new Cup4
cup4.set(56)
cup4.get() is 56	
\end{Verbatim}	
\end{solution}

\end{parts}

\question
\begin{parts}
\part
Use the REPL to create several \texttt{Vector}s, each populated by a different type of data. 
\begin{solution}
\begin{Verbatim}
val v1 = Vector("This", "Is", "A", "Vector", "Of", "Strings")
println(v1)
val v2 = Vector(4, 5, 6, 7, 8)
println(v2)
val v3 = Vector(Range(0, 2), Range(0, 5))
println(v3)

/* You should do this in the REPL
scala> val v1 = Vector("This", "Is", "A", "Vector", "Of", "Strings")
v1: scala.collection.immutable.Vector[java.lang.String] 
	= Vector(This, Is, A, Vector, Of, Strings)

scala> val v2 = Vector(4, 5, 6, 7, 8)
v2: scala.collection.immutable.Vector[Int] = Vector(4, 5, 6, 7, 8)

scala> val v3 = Vector(Range(0, 2), Range(0, 5))
v3: scala.collection.immutable.Vector[scala.collection.immutable.Range] 
	= Vector(Range(0, 1), Range(0, 1, 2, 3, 4))
*/	
\end{Verbatim}	
\end{solution}

\part
Use the REPL to see if you can make a \texttt{Vector} containing other \texttt{Vector}s.
\begin{solution}
\begin{Verbatim}
val v4 = Vector(Vector(0, 1, 2), Vector(3, 4, 5))
println(v4)

/* You should do this in the REPL
scala> val v4 = Vector(Vector(0, 1, 2), Vector(3, 4, 5))
v4: scala.collection.immutable.Vector[scala.collection.immutable.Vector[Int]] 
	= Vector(Vector(0, 1, 2), Vector(3, 4, 5))
*/	
\end{Verbatim}	
\end{solution}

\part
Create a \texttt{Vector} and populate it with words (which are \texttt{String}s). 
Add a \texttt{for} loop that prints each element in the \texttt{Vector}. 
Append to a variable of type \texttt{String} to create a \texttt{sentence}. 
Satisfy the following test:
\begin{Verbatim}
sentence.toString() is "The dog visited the fire station "
\end{Verbatim}

\begin{solution}
\begin{Verbatim}
val v5 = Vector("The", "dog", "visited",
  "the", "fire station")
var sentence: String = ""
for (word <- v5) {
  sentence = sentence + (word + " ")
}

sentence.toString() is
  "The dog visited the fire station "
	
\end{Verbatim}	
\end{solution}

\part
Create and initialise two \texttt{Vector}s, one containing \texttt{Int}s and one containing 
\texttt{Double}s. Call the \texttt{sum}, \texttt{min}, and \texttt{max} 
operations on each one.

\begin{solution}
\begin{Verbatim}
val intVector = Vector(10, 20, 30, 40, 50)
val doubleVector = Vector(10.1D,
  20.2D, 30.3D, 40.4D, 50.5D)
intVector.sum is 150
intVector.min is 10
intVector.max is 50

doubleVector.sum is 151.5D
doubleVector.min is 10.1D
doubleVector.max is 50.5D	
\end{Verbatim}	
\end{solution}

\part
Create two \texttt{Vector}s of \texttt{Int} named 
\texttt{myVector1} and \texttt{myVector2}, each 
initialised to \verb!1, 2, 3, 4, 5, 6!. Use \verb!AtomicTest! to show whether the two are equivalent.

\begin{solution}
\begin{Verbatim}
val myVector1 = Vector(1, 2, 3, 4, 5, 6)
val myVector2 = Vector(1, 2, 3, 4, 5, 6)

myVector1 equals myVector2	
\end{Verbatim}	
\end{solution}

\end{parts}

\question

\begin{parts}
\part
What are the values and types of the following Scala literals?
\begin{subparts}
\subpart \verb!42!
\subpart \verb!true!
\subpart \verb!123L!
\subpart \verb!42.0!
\end{subparts}

\begin{solution}
	\ldots	
\end{solution}

\part
What is the difference between the following literals? 
\begin{subparts}
\subpart \verb!'a'!
\subpart \verb!"a"!
\end{subparts}
What is the type and value of each?

\begin{solution}
	\ldots
\end{solution}

\part
What is the difference between the following expressions? 
\begin{subparts}
\subpart \verb!"Hello world!"!
\subpart \verb!println("Hello world!")!
\end{subparts}
What is the type and value of each?

\begin{solution}
	\ldots
\end{solution}

\part
What is the type and value of the following literal? 
Try writing it using the \textsc{REPL} or in a Scala \emph{worksheet} and see what happens!
\begin{Verbatim}
'Hello world!'
\end{Verbatim}

\begin{solution}
	\ldots
\end{solution}
\end{parts}

\question

\begin{parts}
\part
What are the types and values of the following conditionals?
\begin{subparts}
\subpart
\begin{Verbatim}
val a = 1
val b = 2
if(a > b) "alien" else "predator"
\end{Verbatim}

\begin{solution}
	\ldots
\end{solution}

\subpart
\begin{Verbatim}
val a = 1
val b = 2
if(a > b) "alien" else 2001
\end{Verbatim}

\begin{solution}
	\ldots
\end{solution}

\subpart
\begin{Verbatim}
if(true) "hello"
\end{Verbatim}

\begin{solution}
	\ldots
\end{solution}
\end{subparts}

\part
What is the difference between the following expressions? What are the similarities?
\begin{subparts}
\subpart \verb!1 + 2 + 3!
\subpart \verb!6!
\end{subparts}

\begin{solution}
	\ldots
\end{solution}
\end{parts}

\question
\begin{parts}
\part
Define an object called \verb!calc! 
with a method \verb!square! that accepts a \verb!Double! as an argument and\ldots 
you guessed it\ldots squares its input. 
Add a method called \verb!cube! that cubes its input and calls \verb!square! 
as part of its result calculation.

\begin{solution}
	\ldots
\end{solution}

\part
Copy and paste \verb!calc! from the previous exercise to create a \verb!calc2! 
that is generalised to work with \verb!Int!s 
as well as \verb!Double!s. As you have Java experience, this should be fairly straightforward. 

\begin{solution}
	\dots
\end{solution}
\end{parts}

\question
When entered on the \textsc{REPL}, what does the following program output, 
and what is the type and value of the final
expression? 
\begin{Verbatim}
object argh {
  def a = {
    println("a")
    1
  }

  val b = {
    println("b")
    a + 2
  }

  def c = {
    println("c")
    a
    b + "c"
  }
}

argh.c + argh.b + argh.a
\end{Verbatim}
Think carefully about the types, dependencies, and evaluation behaviour of each field and method.

\begin{solution}
	
\end{solution}

\question
\begin{parts}
\part
Define an object called \verb!Person! that contains fields called \verb!firstName! and 
\verb!lastName!. 

\begin{solution}
	
\end{solution}

\part
Define a second object called \verb!Alien! containing a method called \verb!greet! 
that takes your \verb!Person! as a parameter 
and returns a greeting using their \verb!firstName!.

\begin{solution}
	
\end{solution}
What is the type of the \verb!greet! method? Can we use this method to greet other objects?

\part
Are methods values? Are they expressions? Why might this be the case?

\begin{solution}
	\ldots	
\end{solution}
\end{parts}

\question
Create your own package containing three trivial classes (just define
the classes, don't give them bodies). Now
import one class, two classes, and all classes, and show that you've
successfully imported them in each case.

\begin{solution}
	
\end{solution}

\question
Create a value of type \verb!Range! that goes from \verb!0! to \verb!10! (not including 10).
Satisfy the following tests:
\begin{Verbatim}
val r1 = // fill this in      
r1 is // fill this in 
\end{Verbatim}
Use \verb!Range.inclusive! to solve the problem above. What changed?

\begin{solution}
	
\end{solution}

\question
\begin{parts}
\part
Write a for loop that adds the values 0 through 10 (including 10). Sum
all the values and ensure that it equals 55. Must you use a var instead
of a val ? Why? Satisfy the following test: total is 55

\begin{solution}
	\ldots	
\end{solution}

\part
Write a for loop that adds even numbers between 1 and 10 (including 10).
Sum all the values and ensure that it equals 30. Hint: this conditional
expression determines whether a number is even: if (number \% 2 == 0)
The \% (modulo) operator checks to see if there is a remainder when you
divide number by 2. Satisfy the following:
\begin{verbatim}
`totalEvens is 30`
\end{verbatim}

\begin{solution}
	\ldots
\end{solution}

\part
Write a for loop that adds even numbers between 1 and 10 (including 10)
and odd numbers between 1 and 10. Calculate a sum for the even numbers
and a sum for the odd numbers. Did you write two for loops? If so, try
rewriting this with a single for loop. Satisfy the following tests:
\begin{verbatim}
`evens is 30`
\end{verbatim}

\begin{solution}
	\ldots
\end{solution}
\end{parts}

\question

\begin{parts}
\part
Create a \verb!Vector! and populate it with words (which are \verb!String!s). Add a
for loop that prints each element in the \verb!Vector!. Building on the
previous exercise, append to a variable of type \verb!String! to create a sentence. Satisfy
the following test:
\begin{Verbatim}
sentence.toString() is "The dog visited the firehouse "
\end{Verbatim}
That last space is unexpected. Use \verb!String!'s method \verb!replace! to replace
\verb!firehouse! with \verb!firehouse!! Satisfy the following test:
\begin{Verbatim}
theString is "The dog visited the firehouse!"
\end{Verbatim}

\begin{solution}
	\ldots
\end{solution}

\part
Building on your solution write a \emph{for
loop} that prints each word, reversed. Your output should match:
\begin{Verbatim}
Output: ehT god detisiv eht esuoherif
\end{Verbatim}

\begin{solution}
	\ldots
\end{solution}

\part
Write a \emph{for loop} that prints the words in reverse order
(last word first, etc.). Your output should match:
\begin{Verbatim}
firehouse the visited dog The
\end{Verbatim}

\begin{solution}
	\ldots
\end{solution}
\end{parts}

\question
\begin{parts}
\part
Create and initialise two \verb!Vector!s, one containing \verb!Int!s 
and one
containing \verb!Double!s. Call the \verb!sum!, \verb!min!, and \verb!max! 
operations on each one and
see what happens. 

\begin{solution}
	
\end{solution}

\part
Create a \verb!Vector! containing \verb!String!s and apply the \verb!sum!,
\verb!min!, and \verb!max! operations. 
Explain the results. One of those methods won't work. Why?

\begin{solution}
	
\end{solution}

\part
In \emph{For Loops}, we added the values in a \verb!Range! to get the sum. 
Try calling the \verb!sum! operation on a \verb!Range!. 
Does this do the entire summation in one step?

\begin{solution}
	\ldots
\end{solution}

\part
\verb!List! and \verb!Set! are similar to \verb!Vector!. 
Use the \textsc{REPL} to discover their
operations and compare them to those of \verb!Vector!. 
Create and initialise a
\verb!List! and \verb!Set! with words, then print each one. 
Try the \verb!reverse! and \verb!sorted!
operations and see what happens.

\begin{solution}
	\ldots	
\end{solution}
\end{parts}

\question
Palindromes are words or phrases that read the same forward and
backward. Some examples include ``mum'' and ``dad''. 

Write a method to test words or phrases for palindromes. 

Hint: \verb!String!'s \verb!reverse! method may prove useful here. 

Satisfy the following tests: 
\begin{Verbatim}
isPalindrome(``mum'') 
\end{Verbatim}
is true, and 
\begin{Verbatim}
isPalindrome(``dad'')
\end{Verbatim}
is true, and 
\begin{Verbatim}
isPalindrome(``street'')
\end{Verbatim}
is false.

\begin{solution}
	
\end{solution}

\question
Create a function \verb!forecast! that represents the percentage of cloudiness,
and use it to produce a ``weather forecast''. Strings such as ``Sunny''
(100), ``Mostly Sunny'' (80), ``Partly Sunny'' (50), ``Mostly Cloudy''
(20), and ``Cloudy'' (0). 

For this exercise use pattern matching and only match for the legal
values 100, 80, 50, 20, and 0. Everything else should produce
``Unknown''.

Your function should satisfy the following tests:
\begin{Verbatim}
forecast(100) is "Sunny" 
forecast(80) is "Mostly Sunny" 
forecast(50) is "Partly Sunny" 
forecast(20) is "Mostly Cloudy" 
forecast(0) is "Cloudy" 
forecast(15) is "Unknown"
\end{Verbatim}

\begin{solution}
	
\end{solution}

\question
\begin{parts}
\part
Define a class \verb!SimpleTime! that takes two arguments: an \verb!Int! that
represents hours, and an \verb!Int! that represents minutes. Use \emph{named
arguments} to create a \verb!SimpleTime! object which satisfies the following tests:
\begin{Verbatim}
val t = new SimpleTime(hours=5, minutes=30) 
t.hours is 5
t.minutes is 30
\end{Verbatim}

\begin{solution}
	\ldots	
\end{solution}

\part
Using your answer to the previous question for \verb!SimpleTime!, default the
minutes to 0 so that you don't have to specify them. Your code should satisfy the following tests:
\begin{Verbatim}
val t2 = new SimpleTime2(hours=10)
t2.hours is 10
t2.minutes is 0
\end{Verbatim}

\begin{solution}
	\ldots
\end{solution}
\end{parts}

%\question
%Create a class \verb!Planet! that has, by default, a single moon. 
%The \verb!Planet! class should have a \verb!name! 
%(a \verb!String!) and \verb!description! (a \verb!String!). 
%Use
%\emph{named arguments} to specify the \verb!name! and \verb!description!, 
%and a default for
%the number of moons. Satisfy the following tests:
%\begin{Verbatim}
%val p = new Planet(name = "Mercury", 
%			description = "small and hot planet", moons = 0) 
%p.hasMoon is false
%\end{Verbatim}
%
%\begin{solution}
%	
%\end{solution}

\question
\begin{parts}
\part
Create a \emph{case class} to represent a \verb!Person! in an address book. Complete
the class with a \verb!String! for the \verb!name! and a 
\verb!String! for contact information. Satisfy
the following tests:
\begin{Verbatim}
val p = Person("Jane", "Smile", "jane@smile.com") 
p.first is "Jane" 
p.last is "Smile" 
p.email is "jane@smile.com"
\end{Verbatim}

\begin{solution}
	
\end{solution}

\part
Create some \verb!Person! objects. Put the \verb!Person! objects in a \verb!Vector!. 
Satisfy the following tests:
\begin{Verbatim}
val people = Vector(
    Person("Jane","Smile","jane@smile.com"), 
    Person("Ron","House","ron@house.com"), 
    Person("Sally","Dove","sally@dove.com")
)
people(0) is "Person(Jane,Smile,jane@smile.com)" 
people(1) is "Person(Ron,House,ron@house.com)"
people(2) is "Person(Sally,Dove,sally@dove.com)"
\end{Verbatim}

\begin{solution}
	
\end{solution}

\end{parts}

\question
\begin{parts}
\part
Maps store information using unique keys. An email address can serve as
such a key. Create a class \verb!Name! containing \verb!firstName! and 
\verb!lastName!.

Create a \verb!Map! that associates \verb!emailAddress! (a \verb!String!) with 
\verb!Name!. Satisfy
the following test:
\begin{Verbatim}
val m = Map("sally@taylor.com"-> Name("Sally","Taylor"))
m("sally@taylor.com") is Name("Sally", "Taylor")
\end{Verbatim}

\begin{solution}
	\ldots	
\end{solution}

\part
Augment your solution by adding \verb!Jiminy Cricket! to
the map, where the email address is \verb!jiminy@cricket.com!. 

Satisfy the
following tests: 
\begin{Verbatim}
m2("jiminy@cricket.com") is Name("Jiminy","Cricket") 
m2("sally@taylor.com") is Name("Sally", "Taylor")
\end{Verbatim}

\begin{solution}
	\ldots
\end{solution}
\end{parts}

\question
Sets store distinct values. Create a set for the following languages:
English, French, Spanish, German, and Chinese. Write a Scala class that uses a \verb!Set!
to represent this problem.
\begin{parts}
\part
What happens when you try
to add Turkish? 

\begin{solution}
	
\end{solution}

\part
Try
to add a language that already exists in the set (for example, French).
What happens? 

\begin{solution}
	
\end{solution}

\part
Remove ``Spanish'' from a the set containing English, French,
Spanish, German, and Chinese. 
\part
Remove \verb!jiminy@cricket.com! from a \verb!Map!
containing information for Jiminy Cricket, Mary Smith, and Sally Taylor.
Test your solution.

\begin{solution}
	
\end{solution}

\end{parts}
\end{questions}
\vfill
\section*{Credits}

Thanks to Noel Walsh and Bruce Eckel for the basis of these exercises\ldots
\end{document}